\scanpage{98}
% archive sheet number: 97
quelconque. \inserted{Par expérience personnelle,} Nous pensons cependant\deleted{,} que la compréhension\deleted{, heuristique quoique pour le moment largement heuristique,} obtenue grâce aux arguments de poids et de niveaux, quoiqu'elle soit pour l'instant entièrement heuristique, \deleted{\unsure{est et}} peut être la source de la découverte de propriétés importantes et inattendues pour la cohomologie des variétés algébriques, \deleted{même si} et qu'elle mérite pour cette raison de ne pas être dédaignée. De plus, dans le cas des corps de caractéristique nulle, \deleted{des travaux récents fort \unsure{remar|rem} \unsure{les|las} magnifiques t} le récent et magnifique travail de P. DELIGNE sur la théorie de Hodge\inserted{généralisée} des variétés algébriques quelconques sur le corps $\underline{C}$ des nombres complexes\deleted{,} justifie\inserted{dès maintenant} par voie transcendante le ``yoga'' des poids et niveaux dans le cadre de la cohomologie de De Rham\handwritten{,} (tout en ouvrant à la théorie cohomoligue des variétés algébriques\inserted{en caractéristique nulle} des voies entièrement nouvelles). Il \deleted{est}\inserted{semble} probable d'autre part que pour le cas des \deleted{Pour terminer cette introduction} variétés algé\deleted{\unsure{bu|br}}\handwritten{br}iques sur les corps de caractéristique non nulle, une justification de ce même yoga doive attendre la \deleted{\unsure{résolution|resolution}}\inserted{preuve} des conjectures de Weil (ou mieux, des ``conjectures standart'' de \handwritten{[\quad]}).

Pour \deleted{signaler} terminer cette introduction, signalons que les filtrations par poids et niveaux des espaces de cohomologie $\handwritten{\ell}$-adiques, pour $\handwritten{\ell}$ variable, sont certainement relié\deleted{\unsure{s}}\handwritten{es} entre elles de la façon habituelle, \deleted{donnant}\inserted{et doivent donner} en particulier des espaces vectoriels sur $\underline{Q}_{\handwritten{\ell}}$ \handwritten{d}ont le \deleted{\unsure{ran|rang}} rang est indépendant de $\handwritten{\ell}$. En termes de la théorie conjecturale des motifs (suspendue entre autres aux ``conjectures standart'' déjà mentionnées), ceci devrait s'exprimer par le fait que les structures de poids et de niveau\deleted{se placent}\inserted{des cohomologies $\handwritten{\ell}$-adi\handwritten{ques} se placent déjà} au niveau de\deleted{\unsure{s motifs|motifs} déjà} la \uline{cohomologie} \deleted{des} \uline{motivique}, \deleted{ce qui dans la}\inserted{Cet\deleted{\unsure{te|t}} point de vue est en tous cas en accord avec la} théorie de Deligne, \deleted{\unsure{car}} puisque sa ``filtration par le poids'' de la cohomologie de De Rham provient d'une filtration de la cohomologie à coefficients \uline{rationnels}\handwritten{.}\deleted{(et non seulement}
